\diarytitle{0050}{エアリー方程式のべき級数解と係数の漸化式}

$2$ 階線形常微分方程式 $y'' + P(x)y' + Q(x)y = 0$ において $x=0$ が正則点（$P,Q$ が $x=0$ で正則）であるとき、解は $y = \sum_{n=0}^{\infty} a_n x^n$ の形の収束べき級数として存在する。これを方程式に代入し、各べき $x^m$ の係数を $0$ とおくことで係数 $a_n$ の間の漸化式が得られ、$a_0, a_1$ は任意定数として自由に選べる。

$y = \displaystyle\sum_{n=0}^{\infty} a_n x^n$ より
\[
y'' = \sum_{n=0}^{\infty} (n+2)(n+1)a_{n+2}\,x^{n}, \qquad
xy = \sum_{n=0}^{\infty} a_n x^{n+1} = \sum_{m=1}^{\infty} a_{m-1}\,x^{m}
\]
$y''-xy=0$ の $x^m$ の係数を比較すると、$m=0$ のとき
\[
2\cdot 1\cdot a_2 = 0 \quad\Longrightarrow\quad a_2 = 0
\]
$m\ge 1$ のとき
\[
(m+2)(m+1)a_{m+2} = a_{m-1}
\quad\Longrightarrow\quad
a_{m+2} = \frac{a_{m-1}}{(m+1)(m+2)} \qquad (m \ge 1)
\]
これが漸化式である。これを用いて逐次計算する。
\[
a_3 = \frac{a_0}{3\cdot 2} = \frac{a_0}{6}, \qquad
a_4 = \frac{a_1}{4\cdot 3} = \frac{a_1}{12}, \qquad
a_5 = \frac{a_2}{5\cdot 4} = 0
\]
\[
a_6 = \frac{a_3}{6\cdot 5} = \frac{a_0}{6\cdot 30} = \frac{a_0}{180}, \qquad
a_7 = \frac{a_4}{7\cdot 6} = \frac{a_1}{12\cdot 42} = \frac{a_1}{504}
\]
よって
\[
\boxed{\;
a_2 = 0,\quad a_3=\frac{a_0}{6},\quad a_4=\frac{a_1}{12},\quad a_5=0,\quad a_6=\frac{a_0}{180},\quad a_7=\frac{a_1}{504}
\;}
\]
（検算: 例えば $a_6$ は漸化式で $m=4$ を用いて $a_6 = a_3/(5\cdot 6) = (a_0/6)/30 = a_0/180$ となり一致する。）
